% ============================================================
% RAPPORT DE PROJET DE FIN D'ÉTUDES
% CMAD-Net : Cross-Modal Attention Detection Network
% ============================================================
\documentclass[12pt,a4paper,oneside]{report}

% ── Encodage et langue ──
\usepackage[utf8]{inputenc}
\usepackage[T1]{fontenc}
\usepackage[french]{babel}

% ── Mise en page ──
\usepackage[top=2.5cm, bottom=2.5cm, left=3cm, right=2.5cm]{geometry}
\usepackage{setspace}
\onehalfspacing
\usepackage{fancyhdr}
\usepackage{etoolbox}

% ── Graphiques et couleurs ──
\usepackage{graphicx}
\usepackage{xcolor}
\usepackage{tikz}
\usetikzlibrary{shapes.geometric, arrows.meta, positioning, fit, calc, backgrounds, patterns}

% ── Tableaux ──
\usepackage{booktabs}
\usepackage{tabularx}
\usepackage{array}
\usepackage{longtable}
\usepackage{colortbl}

% ── Code ──
\usepackage{listings}
\usepackage{float}

% ── Liens et PDF ──
\usepackage[hidelinks]{hyperref}
\usepackage{url}

% ── Math ──
\usepackage{amsmath}
\usepackage{amssymb}

% ── Captions ──
\usepackage{caption}
\usepackage{subcaption}

% ── Couleurs personnalisées ──
\definecolor{chaptercolor}{RGB}{25, 80, 140}
\definecolor{sectioncolor}{RGB}{40, 100, 160}
\definecolor{tableheader}{RGB}{220, 235, 250}
\definecolor{lightgray}{RGB}{245, 245, 245}
\definecolor{codegreen}{RGB}{0,120,0}
\definecolor{codegray}{RGB}{100,100,100}
\definecolor{codebg}{RGB}{248,248,248}
\definecolor{accentblue}{RGB}{41, 98, 255}
\definecolor{accentgreen}{RGB}{46, 125, 50}
\definecolor{accentorange}{RGB}{230, 126, 34}
\definecolor{accentred}{RGB}{198, 40, 40}
\definecolor{ganttblue}{RGB}{66,133,244}
\definecolor{ganttgroup}{RGB}{25,80,140}

% ── Style du code ──
\lstset{
  backgroundcolor=\color{codebg},
  basicstyle=\small\ttfamily,
  breaklines=true,
  captionpos=b,
  commentstyle=\color{codegreen},
  keywordstyle=\color{accentblue}\bfseries,
  stringstyle=\color{accentorange},
  numbers=left,
  numberstyle=\tiny\color{codegray},
  numbersep=8pt,
  frame=single,
  rulecolor=\color{codegray!30},
  tabsize=4,
  showstringspaces=false,
  xleftmargin=15pt,
  framexleftmargin=15pt,
}

% ── En-têtes et pieds de page ──
\pagestyle{fancy}
\fancyhf{}
\fancyhead[L]{\small\leftmark}
\fancyhead[R]{\small CMAD-Net}
\fancyfoot[C]{\thepage}
\renewcommand{\headrulewidth}{0.4pt}
\renewcommand{\footrulewidth}{0pt}

% ── Commande checkmark sans pifont ──
\newcommand{\cmark}{\textcolor{accentgreen}{\textbf{\checkmark}}}
\newcommand{\xmark}{\textcolor{accentred}{\textbf{$\times$}}}
\newcommand{\cmadnet}{\textbf{CMAD-Net}}

% ── Profondeur TOC ──
\setcounter{tocdepth}{3}
\setcounter{secnumdepth}{3}

% ============================================================
\begin{document}

% ============================================================
% PAGE DE GARDE
% ============================================================
\begin{titlepage}
\begin{center}

\vspace*{1.5cm}

{\large \textbf{PROJET DE FIN D'ÉTUDES}}

\vspace{0.5cm}
{\large Filière : Génie Informatique / Intelligence Artificielle}

\vspace{2.5cm}

\rule{\textwidth}{1.5pt}
\vspace{0.6cm}
{\Huge \textbf{\color{chaptercolor}CMAD-Net}}\\[0.5cm]
{\LARGE Cross-Modal Attention Detection Network}\\[0.4cm]
{\Large pour la Détection d'Anomalies 3D}
\vspace{0.6cm}
\rule{\textwidth}{1.5pt}

\vspace{2.5cm}

{\large
\begin{tabular}{rl}
\textbf{Réalisé par :} & Souhaib Chettou \\[10pt]
\end{tabular}
}

\vspace{3cm}

{\large \textbf{Année universitaire : 2025 -- 2026}}

\vfill

\end{center}
\end{titlepage}

% ============================================================
% SECTIONS PRÉLIMINAIRES (numérotation romaine)
% ============================================================
\pagenumbering{roman}

% ── Remerciements ──
\chapter*{Remerciements}
\addcontentsline{toc}{chapter}{Remerciements}

Je tiens à exprimer ma profonde gratitude envers toutes les personnes qui ont contribué, de près ou de loin, à la réalisation de ce projet de fin d'études.

Je remercie tout d'abord mon encadrant académique pour ses précieux conseils, son suivi régulier et ses orientations scientifiques qui ont guidé ce travail de recherche.

Je remercie également mon encadrant professionnel pour m'avoir accueilli au sein de son équipe, pour sa confiance et pour les moyens techniques mis à ma disposition.

Mes remerciements vont aussi à l'ensemble des membres du jury pour l'honneur qu'ils me font en évaluant ce travail.

Enfin, je remercie ma famille et mes amis pour leur soutien indéfectible tout au long de mon parcours académique.

% ── Résumé ──
\chapter*{Résumé}
\addcontentsline{toc}{chapter}{Résumé}

Ce projet de fin d'études porte sur la conception et le développement de \textbf{CMAD-Net} (Cross-Modal Attention Detection Network), un réseau de neurones profond pour la détection d'anomalies 3D multi-modale non-supervisée. Le système exploite conjointement les images RGB et les données de profondeur (Depth) via un mécanisme d'attention croisée innovant (CMCA), couplé à un paradigme de Reverse Distillation Teacher-Student. L'architecture intègre un backbone Vision Transformer (ViT-B/16) pré-entraîné avec DINO, un Feature Pyramid Network multi-échelles, et une fusion adaptative par catégorie d'objets. Évalué sur le benchmark MVTec 3D-AD (10 catégories), le modèle atteint un P-AUROC moyen de 0.952 et un AUPRO@30\% de 0.816, démontrant son efficacité pour la localisation précise des anomalies industrielles.

\medskip
\noindent\textbf{Mots-clés :} Détection d'anomalies, Vision par ordinateur, Deep Learning, Vision Transformer, Attention croisée, Multi-modal, 3D, Reverse Distillation.

\vspace{2cm}

\begin{center}
\rule{0.6\textwidth}{0.4pt}
\end{center}

\vspace{0.8cm}

\section*{Abstract}

This final year project focuses on the design and development of \textbf{CMAD-Net} (Cross-Modal Attention Detection Network), a deep neural network for unsupervised multi-modal 3D anomaly detection. The system jointly exploits RGB images and depth data through an innovative cross-modal cross-attention mechanism (CMCA), coupled with a Teacher-Student Reverse Distillation paradigm. The architecture integrates a Vision Transformer backbone (ViT-B/16) pre-trained with DINO, a multi-scale Feature Pyramid Network, and adaptive category-specific score fusion. Evaluated on the MVTec 3D-AD benchmark (10 categories), the model achieves a mean P-AUROC of 0.952 and AUPRO@30\% of 0.816, demonstrating its effectiveness for precise industrial anomaly localization.

\medskip
\noindent\textbf{Keywords:} Anomaly detection, Computer Vision, Deep Learning, Vision Transformer, Cross-attention, Multi-modal, 3D, Reverse Distillation.

% ── Liste des abréviations ──
\chapter*{Liste des Abréviations}
\addcontentsline{toc}{chapter}{Liste des Abréviations}

\begin{longtable}{@{}p{3.5cm}p{10cm}@{}}
\toprule
\textbf{Abréviation} & \textbf{Signification} \\
\midrule
\endhead
AUROC & Area Under the Receiver Operating Characteristic \\
AUPRO & Area Under the Per-Region Overlap \\
BN & Batch Normalization \\
CMAD-Net & Cross-Modal Attention Detection Network \\
CMCA & Cross-Modal Cross-Attention \\
CNN & Convolutional Neural Network \\
DINO & Self-Distillation with NO Labels \\
FPN & Feature Pyramid Network \\
FPR & False Positive Rate \\
GELU & Gaussian Error Linear Unit \\
GPU & Graphics Processing Unit \\
GT & Ground Truth \\
MHA & Multi-Head Attention \\
MSE & Mean Squared Error \\
MVTec 3D-AD & MVTec 3D Anomaly Detection \\
PFE & Projet de Fin d'Études \\
PRO & Per-Region Overlap \\
ReLU & Rectified Linear Unit \\
RGB & Red Green Blue \\
ViT & Vision Transformer \\
\bottomrule
\end{longtable}

% ── Tables des matières, figures, tableaux ──
\tableofcontents
\listoffigures
\addcontentsline{toc}{chapter}{Liste des figures}
\listoftables
\addcontentsline{toc}{chapter}{Liste des tableaux}

% ============================================================
% CORPS DU RAPPORT (numérotation arabe)
% ============================================================
\newpage
\pagenumbering{arabic}

% ============================================================
% INTRODUCTION GÉNÉRALE
% ============================================================
\chapter*{Introduction Générale}
\addcontentsline{toc}{chapter}{Introduction Générale}
\markboth{Introduction Générale}{}

Le contrôle qualité industriel constitue un enjeu stratégique majeur pour les entreprises manufacturières. Dans un contexte de production de masse, la détection automatique des défauts sur les produits est devenue indispensable pour maintenir des standards de qualité élevés tout en optimisant les coûts de production. Les anomalies --- fissures, contaminations, déformations, trous --- peuvent compromettre la fonctionnalité, la sécurité ou l'esthétique des produits, entraînant des pertes financières considérables.

Traditionnellement, l'inspection visuelle est réalisée par des opérateurs humains, une approche coûteuse, subjective et sujette à la fatigue. L'avènement du deep learning et de la vision par ordinateur a ouvert de nouvelles perspectives pour l'automatisation de cette tâche. Cependant, la détection d'anomalies industrielles pose un défi particulier : les anomalies sont, par nature, rares et imprévisibles. C'est pourquoi les approches \textbf{non-supervisées} --- qui n'utilisent que des échantillons normaux pendant l'entraînement --- se sont imposées comme le paradigme dominant.

Un autre défi réside dans la \textbf{multi-modalité}. Les caméras RGB capturent les anomalies visuelles (couleur, texture), mais sont aveugles aux défauts géométriques (bosses, creux). Les capteurs de profondeur capturent la géométrie 3D, mais avec une résolution inférieure. La \textbf{fusion intelligente} de ces deux modalités est essentielle pour une détection complète et robuste.

Dans ce contexte, le présent projet de fin d'études propose \cmadnet{} (Cross-Modal Attention Detection Network), une architecture de réseau de neurones profond qui aborde ces défis à travers trois innovations principales :

\begin{enumerate}
    \item \textbf{L'attention croisée inter-modale (CMCA)} : un mécanisme de fusion profonde qui permet aux features RGB et Depth de s'enrichir mutuellement via un Transformer cross-attention.
    \item \textbf{La Reverse Distillation Teacher-Student} : un paradigme de détection où un student decoder apprend à reconstruire les features d'un teacher sur les échantillons normaux, révélant les anomalies par discordance.
    \item \textbf{Le backbone Vision Transformer (ViT-B/16 DINOv2)} : des représentations auto-supervisées riches extraites par un Vision Transformer pré-entraîné.
\end{enumerate}

Le système est évalué sur le benchmark \textbf{MVTec 3D-AD}, la référence pour la détection d'anomalies 3D industrielles, couvrant 10 catégories d'objets avec des anomalies variées.

Ce rapport est organisé en quatre chapitres :
\begin{itemize}
    \item Le \textbf{Chapitre 1} présente le contexte général du projet, la problématique, les objectifs, la méthodologie et la planification.
    \item Le \textbf{Chapitre 2} dresse un état de l'art des approches de détection d'anomalies, avec une analyse critique et le positionnement de CMAD-Net.
    \item Le \textbf{Chapitre 3} détaille l'analyse des besoins et la conception architecturale du système.
    \item Le \textbf{Chapitre 4} couvre la réalisation, les technologies utilisées et les résultats expérimentaux.
\end{itemize}


% ============================================================
% CHAPITRE 1 : CONTEXTE GÉNÉRAL
% ============================================================
\chapter{Contexte Général du Projet}

\section*{Introduction du chapitre}
\addcontentsline{toc}{section}{Introduction du chapitre}

Ce premier chapitre pose les fondations du projet CMAD-Net. Nous présentons l'organisme d'accueil, le contexte, le périmètre du projet, l'analyse de l'existant, la problématique identifiée et les objectifs fixés. Nous détaillons la méthodologie de développement adoptée et la planification prévisionnelle sous forme de diagramme de Gantt.

\section{Organisme d'accueil}

Le projet CMAD-Net a été réalisé dans le cadre d'un projet de fin d'études en Génie Informatique, spécialité Intelligence Artificielle. Le travail s'inscrit dans le domaine de la recherche appliquée en vision par ordinateur et deep learning, avec un focus sur la détection d'anomalies industrielles 3D.

L'environnement de travail dispose des ressources suivantes :
\begin{itemize}
    \item \textbf{Infrastructure de calcul} : Machines locales et plateforme cloud Kaggle (GPU NVIDIA)
    \item \textbf{Données} : Dataset MVTec 3D-AD (benchmark académique de référence)
    \item \textbf{Outils} : Environnement Python avec PyTorch, bibliothèques de vision par ordinateur
\end{itemize}

\section{Contexte du projet}

\subsection{Périmètre et description du projet}

Le projet CMAD-Net consiste en la \textbf{conception, le développement et l'évaluation} d'un système de détection d'anomalies 3D basé sur le deep learning multi-modal. Le périmètre couvre :

\begin{itemize}
    \item \textbf{L'ingénierie des données} : Pipeline de chargement et pré-traitement des images RGB et des données de profondeur 3D (conversion depth $\rightarrow$ normal maps)
    \item \textbf{L'architecture du modèle} : Conception d'un réseau intégrant un Vision Transformer, un FPN, un mécanisme d'attention croisée et un décodeur student
    \item \textbf{L'entraînement} : Pipeline en 2 phases (Reverse Distillation + Calibration)
    \item \textbf{L'évaluation} : Mesure des performances sur 10 catégories avec un ensemble complet de métriques
\end{itemize}

\subsection{Nature du projet}

Ce projet est un \textbf{sujet de recherche appliquée} combinant recherche bibliographique, innovation architecturale, développement logiciel et évaluation expérimentale. Il s'agit d'une \textbf{évolution} du modèle TRD (Teacher Reverse Distillation) existant, avec des améliorations significatives :

\begin{itemize}
    \item Remplacement du backbone CNN (WideResNet-50) par un Vision Transformer (ViT-B/16 DINOv2)
    \item Introduction d'un mécanisme d'attention croisée (CMCA) remplaçant les Crossmodal Tuners CNN
    \item Adoption du paradigme Reverse Distillation en remplacement du Memory Bank
    \item Post-traitement non-paramétrique remplaçant le FP Filter paramétrique
\end{itemize}

\subsection{L'existant}

Avant CMAD-Net, plusieurs solutions existaient pour la détection d'anomalies 3D :

\textbf{a) TRD (Teacher Reverse Distillation)} --- Backbone WideResNet-50, fusion par Crossmodal Tuners CNN, détection par Memory Bank, post-traitement par FP Filter paramétrique. \textit{Limitation} : fusion tardive et superficielle ; le FP Filter peut supprimer les vraies anomalies.

\textbf{b) PatchCore} --- Memory Bank de patches locaux avec recherche de plus proches voisins. \textit{Limitation} : coût mémoire et temps d'inférence élevés ; principalement 2D.

\textbf{c) RD4AD} --- Teacher-Student avec décodeur inversé. \textit{Limitation} : unimodal (RGB uniquement).

\textbf{d) AST} --- Student-Teacher avec architectures asymétriques. \textit{Limitation} : pas de mécanisme cross-modal explicite.

\subsection{Problématique}

L'analyse de l'existant révèle la problématique suivante :

\begin{quote}
\textit{\textbf{Problématique principale :} Comment concevoir un système de détection d'anomalies 3D capable de fusionner \textbf{profondément} et \textbf{dynamiquement} les informations visuelles (RGB) et géométriques (Depth), tout en maintenant une détection précise au niveau pixel ?}
\end{quote}

Les sous-problèmes identifiés sont :
\begin{enumerate}
    \item \textbf{Fusion superficielle} : Les méthodes existantes combinent les modalités tardivement
    \item \textbf{Représentations limitées} : Les backbones CNN capturent des features locales, manquant le contexte global
    \item \textbf{Détection mono-échelle} : Beaucoup de méthodes n'exploitent que 3 niveaux de features
    \item \textbf{Post-traitement risqué} : Les modules paramétriques peuvent supprimer les vraies anomalies
    \item \textbf{Fusion rigide} : L'importance des modalités varie selon le type d'objet
\end{enumerate}

\subsection{Objectifs et solution proposée}

Le Tableau~\ref{tab:objectifs} présente les objectifs du projet.

\begin{table}[H]
\centering
\caption{Objectifs du projet CMAD-Net}
\label{tab:objectifs}
\begin{tabularx}{\textwidth}{c X c}
\toprule
\textbf{N\textsuperscript{o}} & \textbf{Objectif} & \textbf{Priorité} \\
\midrule
\rowcolor{tableheader} O1 & Concevoir une architecture multi-modale avec fusion profonde par attention croisée & Haute \\
O2 & Implémenter le paradigme Reverse Distillation pour la détection d'anomalies & Haute \\
\rowcolor{tableheader} O3 & Intégrer un backbone Vision Transformer (ViT-B/16 DINOv2) & Haute \\
O4 & Développer un FPN multi-échelles à 4 niveaux & Moyenne \\
\rowcolor{tableheader} O5 & Implémenter une fusion adaptative par catégorie & Moyenne \\
O6 & Évaluer sur les 10 catégories de MVTec 3D-AD & Haute \\
\rowcolor{tableheader} O7 & Atteindre des performances compétitives (P-AUROC $>$ 0.95) & Haute \\
\bottomrule
\end{tabularx}
\end{table}

Le Tableau~\ref{tab:solutions} résume comment CMAD-Net adresse chaque sous-problème.

\begin{table}[H]
\centering
\caption{Solutions proposées par CMAD-Net}
\label{tab:solutions}
\begin{tabularx}{\textwidth}{X X}
\toprule
\textbf{Sous-problème} & \textbf{Solution CMAD-Net} \\
\midrule
\rowcolor{tableheader} Fusion superficielle & CMCA : Attention croisée symétrique à chaque niveau FPN \\
Représentations limitées & ViT-B/16 DINOv2 : Backbone Transformer avec attention globale \\
\rowcolor{tableheader} Détection mono-échelle & FPN à 4 niveaux avec pathway top-down \\
Post-traitement risqué & Post-traitement non-paramétrique (Gaussian smoothing) \\
\rowcolor{tableheader} Fusion rigide & Adaptive Fusion : Poids Softmax appris par catégorie \\
\bottomrule
\end{tabularx}
\end{table}

\subsection{Mission dans le projet}

La mission couvre l'intégralité du cycle de vie :
\begin{enumerate}
    \item \textbf{Recherche bibliographique} : Étude de l'état de l'art
    \item \textbf{Conception architecturale} : Design des 6 modules
    \item \textbf{Développement} : Implémentation en Python/PyTorch
    \item \textbf{Entraînement et optimisation} : Entraînement sur MVTec 3D-AD
    \item \textbf{Évaluation et analyse} : Mesure des performances
    \item \textbf{Rédaction} : Documentation et rapport de PFE
\end{enumerate}

\subsection{Livrables}

Le Tableau~\ref{tab:livrables} présente les livrables du projet.

\begin{table}[H]
\centering
\caption{Livrables du projet}
\label{tab:livrables}
\begin{tabularx}{\textwidth}{c X X c}
\toprule
\textbf{N\textsuperscript{o}} & \textbf{Livrable} & \textbf{Description} & \textbf{Format} \\
\midrule
\rowcolor{tableheader} L1 & Code source CMAD-Net & Architecture complète (7 modules) & Python \\
L2 & Script d'entraînement & Pipeline en 2 phases avec logging & Python \\
\rowcolor{tableheader} L3 & Module d'évaluation & 10 métriques standardisées & Python \\
L4 & Dataset pipeline & Chargement et pré-traitement & Python \\
\rowcolor{tableheader} L5 & Checkpoints & Poids pour 10 catégories & .pth \\
L6 & Logs d'entraînement & Courbes de loss et métriques & Texte \\
\rowcolor{tableheader} L7 & Documentation & Architecture et code & Markdown \\
L8 & Rapport de PFE & Le présent document & PDF \\
\bottomrule
\end{tabularx}
\end{table}

\subsection{Méthodologie de développement}

Nous avons adopté la méthodologie \textbf{Scrum adaptée} (Agile), avec des sprints de 2 semaines. Ce choix se justifie par la nature exploratoire du projet et la nécessité d'adaptabilité.

Le Tableau~\ref{tab:sprints} détaille la planification des sprints.

\begin{table}[H]
\centering
\caption{Planification des sprints Scrum}
\label{tab:sprints}
\begin{tabularx}{\textwidth}{c c X X}
\toprule
\textbf{Sprint} & \textbf{Durée} & \textbf{Objectifs} & \textbf{Livrable} \\
\midrule
\rowcolor{tableheader} Sprint 0 & Sem. 1--2 & Recherche biblio, installation env. & État de l'art \\
Sprint 1 & Sem. 3--4 & ViT Encoder + FPN & Backbone \\
\rowcolor{tableheader} Sprint 2 & Sem. 5--6 & CMCA + Student Decoder & Pipeline T-S \\
Sprint 3 & Sem. 7--8 & Intégration + Fusion + Loss & Modèle complet \\
\rowcolor{tableheader} Sprint 4 & Sem. 9--10 & Dataset + entraînement & 1\textsuperscript{er} entraînement \\
Sprint 5 & Sem. 11--12 & Optimisation + évaluation & 10 catégories \\
\rowcolor{tableheader} Sprint 6 & Sem. 13--14 & Analyse + rédaction & Rapport final \\
\bottomrule
\end{tabularx}
\end{table}

\subsection{Planification prévisionnelle}

La Figure~\ref{fig:gantt} présente le diagramme de Gantt du projet, couvrant 14 semaines.

\begin{figure}[H]
\centering
\resizebox{\textwidth}{!}{%
\begin{tikzpicture}[
    x=1.2cm,
    y=-0.8cm,
]
% ── Grille ──
\foreach \i in {1,...,14} {
    \draw[gray!20] (\i, 0) -- (\i, -16);
    \node[above, font=\small\bfseries] at (\i+0.5, 0) {S\i};
}
\foreach \i in {0,...,16} {
    \draw[gray!20] (1, -\i) -- (15, -\i);
}

% ── Titre ──
\node[above, font=\large\bfseries\color{chaptercolor}] at (8, 0.5) {Diagramme de Gantt --- Projet CMAD-Net};

% ── Sprints (groupes) ──
\fill[ganttgroup, rounded corners=2pt] (1, -0.5) rectangle (3, -1) node[midway, white, font=\small\bfseries] {Sprint 0 --- Recherche};
\fill[ganttblue!60, rounded corners=2pt] (1, -1.3) rectangle (2.8, -1.7) node[midway, font=\scriptsize] {État de l'art};
\fill[ganttblue!60, rounded corners=2pt] (1, -1.9) rectangle (1.9, -2.3) node[midway, font=\scriptsize] {Env. setup};

\fill[ganttgroup, rounded corners=2pt] (3, -2.8) rectangle (5, -3.3) node[midway, white, font=\small\bfseries] {Sprint 1 --- Backbone};
\fill[ganttblue!60, rounded corners=2pt] (3, -3.6) rectangle (4, -4) node[midway, font=\scriptsize] {ViT Encoder};
\fill[ganttblue!60, rounded corners=2pt] (4, -3.6) rectangle (5, -4) node[midway, font=\scriptsize] {FPN};

\fill[ganttgroup, rounded corners=2pt] (5, -4.5) rectangle (7, -5) node[midway, white, font=\small\bfseries] {Sprint 2 --- Modules clés};
\fill[ganttblue!60, rounded corners=2pt] (5, -5.3) rectangle (6.2, -5.7) node[midway, font=\scriptsize] {CMCA};
\fill[ganttblue!60, rounded corners=2pt] (6, -5.3) rectangle (7, -5.7) node[midway, font=\scriptsize] {Student Dec.};

\fill[ganttgroup, rounded corners=2pt] (7, -6.2) rectangle (9, -6.7) node[midway, white, font=\small\bfseries] {Sprint 3 --- Intégration};
\fill[ganttblue!60, rounded corners=2pt] (7, -7) rectangle (9, -7.4) node[midway, font=\scriptsize] {CMADNet + Fusion + Loss};

\fill[ganttgroup, rounded corners=2pt] (9, -7.9) rectangle (11, -8.4) node[midway, white, font=\small\bfseries] {Sprint 4 --- Dataset \& Training};
\fill[ganttblue!60, rounded corners=2pt] (9, -8.7) rectangle (10, -9.1) node[midway, font=\scriptsize] {Dataset};
\fill[ganttblue!60, rounded corners=2pt] (10, -8.7) rectangle (11, -9.1) node[midway, font=\scriptsize] {Training};

\fill[ganttgroup, rounded corners=2pt] (11, -9.6) rectangle (13, -10.1) node[midway, white, font=\small\bfseries] {Sprint 5 --- Évaluation};
\fill[ganttblue!60, rounded corners=2pt] (11, -10.4) rectangle (12.2, -10.8) node[midway, font=\scriptsize] {10 catégories};
\fill[ganttblue!60, rounded corners=2pt] (12, -10.4) rectangle (13, -10.8) node[midway, font=\scriptsize] {Métriques};

\fill[ganttgroup, rounded corners=2pt] (13, -11.3) rectangle (15, -11.8) node[midway, white, font=\small\bfseries] {Sprint 6 --- Rédaction};
\fill[ganttblue!60, rounded corners=2pt] (13, -12.1) rectangle (15, -12.5) node[midway, font=\scriptsize] {Rapport PFE + Soutenance};

% ── Flèches de dépendance ──
\draw[-{Stealth}, thick, gray] (2.8, -1.5) -- (3, -3.8);
\draw[-{Stealth}, thick, gray] (5, -3.8) -- (5, -5.5);
\draw[-{Stealth}, thick, gray] (7, -5.5) -- (7, -7.2);
\draw[-{Stealth}, thick, gray] (9, -7.2) -- (9, -8.9);
\draw[-{Stealth}, thick, gray] (11, -8.9) -- (11, -10.6);
\draw[-{Stealth}, thick, gray] (13, -10.6) -- (13, -12.3);

\end{tikzpicture}
}%
\caption{Diagramme de Gantt prévisionnel du projet CMAD-Net}
\label{fig:gantt}
\end{figure}

La Figure~\ref{fig:gantt} montre la répartition temporelle sur 14 semaines. Les dépendances entre sprints sont indiquées par les flèches. Chaque sprint produit un incrément fonctionnel, conformément à la méthodologie Scrum.

\section*{Conclusion du chapitre}
\addcontentsline{toc}{section}{Conclusion du chapitre}

Ce premier chapitre a posé les bases du projet CMAD-Net. Nous avons présenté le contexte, le périmètre, l'analyse de l'existant, la problématique de la fusion multi-modale profonde et les objectifs. La méthodologie Scrum et le diagramme de Gantt illustrent la planification. Le chapitre suivant dresse un état de l'art approfondi.


% ============================================================
% CHAPITRE 2 : ÉTAT DE L'ART
% ============================================================
\chapter{État de l'Art}

\section*{Introduction du chapitre}
\addcontentsline{toc}{section}{Introduction du chapitre}

Ce chapitre dresse un panorama des approches de détection d'anomalies en vision par ordinateur, avec un focus sur les méthodes multi-modales 3D. Nous définissons les concepts clés, analysons critiquement les solutions existantes, et positionnons CMAD-Net.

\section{Concepts clés}

\subsection{Détection d'anomalies industrielles}

La détection d'anomalies industrielles consiste à identifier automatiquement les défauts sur des produits manufacturés. Ce problème se décompose en :
\begin{itemize}
    \item \textbf{Classification d'images} (Image-level) : binaire normal/anormal
    \item \textbf{Segmentation d'anomalies} (Pixel-level) : localisation précise des régions anomales
\end{itemize}

Les approches non-supervisées (entraînement sur échantillons normaux uniquement) dominent car les anomalies sont rares et imprévisibles.

\subsection{Vision Transformer (ViT)}

Introduit par Dosovitskiy et al.~\cite{dosovitskiy2021}, le ViT découpe l'image en patches de 16$\times$16 pixels, les projette en embeddings, et les traite comme une séquence de tokens via des blocs Transformer. Les avantages pour la détection d'anomalies : champ réceptif global, features multi-niveaux et pré-entraînement auto-supervisé DINO \cite{caron2021}.

\subsection{Attention croisée (Cross-Attention)}

L'attention croisée est une variante où Q provient d'une source et K/V d'une autre :
\begin{equation}
\text{CrossAttention}(Q_A, K_B, V_B) = \text{softmax}\left(\frac{Q_A K_B^T}{\sqrt{d_k}}\right) V_B
\end{equation}

Ce mécanisme permet une interaction dynamique entre deux modalités.

\subsection{Reverse Distillation}

Le paradigme de Reverse Distillation \cite{deng2022} inverse le knowledge distillation :
\begin{itemize}
    \item \textbf{Teacher} : Encodeur pré-entraîné (gelé) qui extrait des features riches
    \item \textbf{Student} : Décodeur reconstruisant les features multi-échelles du teacher
    \item \textbf{Détection} : Discordance $1 - \cos(\text{teacher}, \text{student})$ révèle les anomalies
\end{itemize}

\subsection{Feature Pyramid Network (FPN)}

Le FPN \cite{lin2017} est une architecture multi-échelles avec pathway top-down et connexions latérales, permettant une détection efficace à toutes les échelles spatiales.

\subsection{Benchmark MVTec 3D-AD}

Le dataset MVTec 3D-AD \cite{bergmann2022} contient 10 catégories (bagel, cable\_gland, carrot, cookie, dowel, foam, peach, potato, rope, tire) avec images RGB, nuages de points 3D et masques ground truth.

\section{Analyse critique des solutions existantes}

\subsection{Approches par reconstruction}

\textbf{Principe :} Un modèle (auto-encodeur, GAN) reconstruit les images normales. Les anomalies produisent une erreur élevée.

\textit{Forces :} Simplicité. \textit{Faiblesses :} Reconstruction possible des anomalies ; instabilité des GANs ; unimodales.

\subsection{Approches par Memory Bank}

\textbf{Principe :} Stockage des embeddings de patches normaux. Distance au plus proche voisin = score d'anomalie. Ex. : PatchCore \cite{roth2022}.

\textit{Forces :} Pas d'entraînement ; performances élevées en 2D. \textit{Faiblesses :} Coût mémoire $O(N \times D)$ ; inférence $O(N)$.

\subsection{Approches Teacher-Student}

\textbf{Principe :} Le student reproduit les sorties du teacher sur les données normales. Ex. : STPM \cite{wang2021}.

\textit{Forces :} Inférence rapide. \textit{Faiblesses :} Risque de copie ; unimodales.

\subsection{Approches Reverse Distillation}

\textbf{Principe :} Le student est un \textbf{décodeur} reconstruisant les features du teacher. Ex. : RD4AD \cite{deng2022}.

\textit{Forces :} Bottleneck empêche la copie ; multi-échelles naturel. \textit{Faiblesses :} Unimodal ; pas de fusion cross-modale.

\subsection{Approches multi-modales 3D}

\textit{Forces :} Détection d'anomalies visuelles ET géométriques. \textit{Faiblesses :} Fusion souvent tardive ; pas d'attention croisée profonde.

\subsection{Tableau comparatif}

Le Tableau~\ref{tab:comparaison} compare les principales approches. CMAD-Net est la seule combinant fusion profonde, ViT, attention croisée et post-traitement non-paramétrique.

\begin{table}[H]
\centering
\caption{Comparaison des approches de détection d'anomalies}
\label{tab:comparaison}
\small
\begin{tabularx}{\textwidth}{X c c c c c c}
\toprule
\textbf{Critère} & \rotatebox{65}{\textbf{Reconstr.}} & \rotatebox{65}{\textbf{MemBank}} & \rotatebox{65}{\textbf{T-S}} & \rotatebox{65}{\textbf{Rev.Dist.}} & \rotatebox{65}{\textbf{TRD}} & \rotatebox{65}{\textbf{CMAD-Net}} \\
\midrule
\rowcolor{tableheader} Multi-modal & \xmark & \xmark & \xmark & \xmark & Tardive & \textbf{Profonde} \\
Multi-échelles & Faible & Oui & Oui & Oui & Oui & \textbf{4 niveaux} \\
\rowcolor{tableheader} Coût mémoire & Mod. & \textbf{Élevé} & Faible & Faible & Mod. & \textbf{Faible} \\
Backbone & CNN & CNN & CNN & WRN-50 & WRN-50 & \textbf{ViT DINOv2} \\
\rowcolor{tableheader} Cross-Attn & \xmark & \xmark & \xmark & \xmark & \xmark & \textbf{CMCA} \\
Post-trait. & -- & -- & -- & -- & FP Filter & \textbf{Non-param.} \\
\bottomrule
\end{tabularx}
\end{table}

Le Tableau~\ref{tab:backbones} compare les backbones.

\begin{table}[H]
\centering
\caption{Comparaison des backbones pour la détection d'anomalies}
\label{tab:backbones}
\begin{tabularx}{\textwidth}{X c c c c}
\toprule
\textbf{Backbone} & \textbf{Type} & \textbf{Pré-entraîn.} & \textbf{Champ réc.} & \textbf{Params} \\
\midrule
\rowcolor{tableheader} ResNet-18 & CNN & ImageNet sup. & Local & 11M \\
WideResNet-50 & CNN & ImageNet sup. & Local & 68M \\
\rowcolor{tableheader} EfficientNet-B4 & CNN & ImageNet sup. & Local & 19M \\
\textbf{ViT-B/16 DINOv2} & \textbf{Transf.} & \textbf{DINO auto-sup.} & \textbf{Global} & \textbf{86M} \\
\bottomrule
\end{tabularx}
\end{table}

\section{Synthèse et positionnement}

L'analyse critique révèle cinq lacunes principales :
\begin{enumerate}
    \item \textbf{Absence de fusion profonde} entre les modalités RGB et Depth
    \item \textbf{Sous-exploitation des Vision Transformers}
    \item \textbf{Post-traitement risqué} (FP Filter paramétrique)
    \item \textbf{Manque d'adaptabilité} des poids de fusion par catégorie
    \item \textbf{Détection limitée à 3 niveaux} de features
\end{enumerate}

CMAD-Net comble ces lacunes en proposant respectivement : le CMCA, le ViT-B/16 DINOv2, le post-traitement non-paramétrique, l'Adaptive Fusion, et le FPN à 4 niveaux.

\section*{Conclusion du chapitre}
\addcontentsline{toc}{section}{Conclusion du chapitre}

Ce chapitre a présenté un panorama complet des approches de détection d'anomalies. CMAD-Net se positionne comme une solution innovante combinant un backbone ViT, un mécanisme CMCA, le paradigme Reverse Distillation et une fusion adaptative. Le chapitre suivant détaille la conception architecturale.


% ============================================================
% CHAPITRE 3 : ANALYSE ET CONCEPTION
% ============================================================
\chapter{Analyse et Conception}

\section*{Introduction du chapitre}
\addcontentsline{toc}{section}{Introduction du chapitre}

Ce chapitre constitue le c\oe{}ur technique du projet. Nous présentons l'analyse des besoins, puis la conception architecturale détaillée de chaque module de CMAD-Net avec les diagrammes UML et l'architecture technique.

\section{Analyse des besoins}

\subsection{Besoins fonctionnels}

Le Tableau~\ref{tab:besoins_fonc} liste les besoins fonctionnels.

\begin{table}[H]
\centering
\caption{Besoins fonctionnels du système}
\label{tab:besoins_fonc}
\small
\begin{tabularx}{\textwidth}{c X c}
\toprule
\textbf{ID} & \textbf{Besoin} & \textbf{Priorité} \\
\midrule
\rowcolor{tableheader} BF1 & Charger et pré-traiter les images RGB et données de profondeur 3D & Haute \\
BF2 & Extraire des features multi-échelles via un backbone ViT & Haute \\
\rowcolor{tableheader} BF3 & Projeter les features par un FPN multi-niveaux & Haute \\
BF4 & Fusionner les features RGB et Depth par attention croisée & Haute \\
\rowcolor{tableheader} BF5 & Reconstruire les features teacher via un student decoder & Haute \\
BF6 & Calculer les cartes d'anomalies par discordance teacher-student & Haute \\
\rowcolor{tableheader} BF7 & Fusionner les cartes RGB et Depth avec des poids adaptifs & Moyenne \\
BF8 & Évaluer avec les métriques AUROC, AUPRO, AP, F1 & Haute \\
\rowcolor{tableheader} BF9 & Sauvegarder et charger les poids du modèle & Moyenne \\
BF10 & Supporter l'entraînement sur GPU et CPU & Moyenne \\
\bottomrule
\end{tabularx}
\end{table}

\subsection{Besoins non-fonctionnels}

\begin{itemize}
    \item \textbf{Performance} : P-AUROC $>$ 0.95 sur MVTec 3D-AD
    \item \textbf{Modularité} : Architecture en modules indépendants et réutilisables
    \item \textbf{Reproductibilité} : Seed fixe, déterminisme
    \item \textbf{Documentation} : Code documenté avec docstrings
    \item \textbf{Extensibilité} : Extensible à d'autres backbones ou modalités
    \item \textbf{Compatibilité} : Machine locale et plateforme Kaggle
\end{itemize}

\section{Conception architecturale}

\subsection{Architecture globale de CMAD-Net}

La Figure~\ref{fig:archi_globale} présente l'architecture globale de CMAD-Net.

\begin{figure}[H]
\centering
\begin{tikzpicture}[
    node distance=0.7cm and 1cm,
    block/.style={rectangle, draw, rounded corners=3pt, minimum height=0.9cm, minimum width=2.8cm, text centered, font=\small},
    arrow/.style={-{Stealth[length=2.5mm]}, thick},
]

% Entrées
\node[block, fill=green!20] (rgb) {Image RGB\\{\scriptsize(B, 3, 224, 224)}};
\node[block, fill=blue!20, right=2.5cm of rgb] (depth) {Normal Map\\{\scriptsize(B, 3, 224, 224)}};

% ViT
\node[block, fill=gray!15, below=0.9cm of $(rgb)!0.5!(depth)$, minimum width=6cm] (vit) {\textbf{ViT-B/16 DINOv2} (Gelé) --- Couches 3, 6, 9, 12};

% FPN
\node[block, fill=indigo!8, below left=0.9cm and 0.2cm of vit] (fpn_rgb) {FPN RGB\\{\scriptsize 768$\to$256}};
\node[block, fill=indigo!8, below right=0.9cm and 0.2cm of vit] (fpn_depth) {FPN Depth\\{\scriptsize 768$\to$256}};

% CMCA
\node[block, fill=purple!15, below=0.9cm of $(fpn_rgb)!0.5!(fpn_depth)$, minimum width=5cm] (cmca) {\textbf{CMCA} --- 4 niveaux $\times$ 8 têtes --- Attention croisée};

% Student
\node[block, fill=yellow!12, below left=0.9cm and 0.2cm of cmca] (stu_rgb) {Student RGB\\{\scriptsize Bottleneck + 4 Blocks}};
\node[block, fill=yellow!12, below right=0.9cm and 0.2cm of cmca] (stu_depth) {Student Depth\\{\scriptsize Bottleneck + 4 Blocks}};

% Fusion
\node[block, fill=orange!15, below=0.9cm of $(stu_rgb)!0.5!(stu_depth)$, minimum width=5.5cm] (fusion) {\textbf{Adaptive Fusion} --- $w_{\text{rgb}} \cdot \text{map}_{\text{rgb}} + w_{\text{depth}} \cdot \text{map}_{\text{depth}}$};

% Sortie
\node[block, fill=red!15, below=0.7cm of fusion] (output) {\textbf{Anomaly Map} {\scriptsize(B, 1, 224, 224)}};

% Flèches
\draw[arrow] (rgb) -- (vit);
\draw[arrow] (depth) -- (vit);
\draw[arrow] (vit) -- (fpn_rgb);
\draw[arrow] (vit) -- (fpn_depth);
\draw[arrow] (fpn_rgb) -- (cmca);
\draw[arrow] (fpn_depth) -- (cmca);
\draw[arrow] (cmca) -- node[left, font=\tiny] {P4} (stu_rgb);
\draw[arrow] (cmca) -- node[right, font=\tiny] {P4} (stu_depth);
\draw[arrow] (stu_rgb) -- (fusion);
\draw[arrow] (stu_depth) -- (fusion);
\draw[arrow] (fusion) -- (output);

% Teacher feats vers fusion (pointillé)
\draw[arrow, dashed, gray] (cmca.south west) -- ++(0,-0.3) -| ([xshift=-5mm]fusion.north west) node[near start, above, font=\tiny, gray] {teacher feats};

\end{tikzpicture}
\caption{Architecture globale de CMAD-Net}
\label{fig:archi_globale}
\end{figure}

La Figure~\ref{fig:archi_globale} montre le modèle organisé en trois chemins : le chemin Teacher (ViT $\to$ FPN $\to$ CMCA), le chemin Student (Student Decoders), et la couche de détection (Cosine Distance $\to$ Fusion). Le backbone ViT est partagé entre les deux modalités et gelé.

\subsection{Diagramme de classes UML}

La Figure~\ref{fig:classes_uml} présente le diagramme de classes. La classe centrale \texttt{CMADNet} agrège tous les modules par composition.

\begin{figure}[H]
\centering
\begin{tikzpicture}[
    umlclass/.style={rectangle, draw, rounded corners=1pt, font=\small, minimum width=3.5cm, text width=3.3cm, align=left},
    comp/.style={-{Diamond[open, length=3mm, width=2.5mm]}, thick},
    node distance=0.5cm,
]

% CMADNet
\node[umlclass, fill=chaptercolor!8] (cmad) {
\textbf{CMADNet}\\[3pt]
{\scriptsize +encoder: ViTDINOv2Encoder}\\
{\scriptsize +fpn\_rgb: MultiScaleFPN}\\
{\scriptsize +fpn\_depth: MultiScaleFPN}\\
{\scriptsize +cmca: CMCA}\\
{\scriptsize +student\_rgb: StudentDecoder}\\
{\scriptsize +student\_depth: StudentDecoder}\\
{\scriptsize +fusion: AdaptiveFusion}\\[3pt]
{\scriptsize +forward(rgb, depth)}\\
{\scriptsize +get\_trainable\_params()}\\
{\scriptsize +compute\_anomaly\_maps()}
};

% ViT
\node[umlclass, fill=green!5, right=2cm of cmad, yshift=2cm] (vit) {
\textbf{ViTDINOv2Encoder}\\[3pt]
{\scriptsize -model: timm.ViT}\\
{\scriptsize -embed\_dim: 768}\\
{\scriptsize +forward(x) $\to$ list}
};

% FPN
\node[umlclass, fill=blue!5, right=2cm of cmad, yshift=0cm] (fpn) {
\textbf{MultiScaleFPN}\\[3pt]
{\scriptsize -lateral\_convs}\\
{\scriptsize -smooth\_convs}\\
{\scriptsize +forward(feats)}
};

% CMCA
\node[umlclass, fill=purple!5, right=2cm of cmad, yshift=-2cm] (cmca_c) {
\textbf{CMCA}\\[3pt]
{\scriptsize -levels: ModuleList}\\
{\scriptsize +forward(f2d, f3d)}
};

% Student
\node[umlclass, fill=yellow!5, below=1.5cm of cmad, xshift=-2cm] (stu) {
\textbf{StudentDecoder}\\[3pt]
{\scriptsize -bottleneck}\\
{\scriptsize -decoder\_blocks}\\
{\scriptsize +forward(feat)}
};

% Fusion
\node[umlclass, fill=orange!5, below=1.5cm of cmad, xshift=2cm] (fus) {
\textbf{AdaptiveFusion}\\[3pt]
{\scriptsize -weights: (10, 2)}\\
{\scriptsize +forward(rgb, depth, id)}
};

% Liens
\draw[comp] (cmad.east) -- ++(0.5,0) |- (vit.west);
\draw[comp] (cmad.east) -- (fpn.west);
\draw[comp] (cmad.east) -- ++(0.5,0) |- (cmca_c.west);
\draw[comp] (cmad.south) -- ++(0,-0.3) -| (stu.north);
\draw[comp] (cmad.south) -- ++(0,-0.3) -| (fus.north);

\end{tikzpicture}
\caption{Diagramme de classes UML de CMAD-Net}
\label{fig:classes_uml}
\end{figure}

\subsection{Spécifications des modules}

Le Tableau~\ref{tab:specs_modules} résume les spécifications de chaque module.

\begin{table}[H]
\centering
\caption{Spécifications des modules de CMAD-Net}
\label{tab:specs_modules}
\begin{tabularx}{\textwidth}{X X X c}
\toprule
\textbf{Module} & \textbf{Entrée} & \textbf{Sortie} & \textbf{Params} \\
\midrule
\rowcolor{tableheader} ViTDINOv2Encoder & (B,3,224,224) & [(B,768,14,14)]$\times$4 & 0 (gelé) \\
MultiScaleFPN ($\times$2) & [(B,768,14,14)]$\times$4 & [(B,256,14,14)]$\times$4 & $\sim$3.2M \\
\rowcolor{tableheader} CMCA & [(B,256,14,14)]$\times$4 & idem ($\times$2 mod.) & $\sim$6.3M \\
StudentDecoder ($\times$2) & (B,256,14,14) & [(B,256,14,14)]$\times$4 & $\sim$2.4M \\
\rowcolor{tableheader} AdaptiveFusion & (B,1,H,W)$\times$2 & (B,1,H,W) & 20 \\
\textbf{Total entraînable} & & & \textbf{$\sim$17.5M} \\
\bottomrule
\end{tabularx}
\end{table}

\section{Conception détaillée des modules}

\subsection{ViT-B/16 DINOv2 Encoder}

Le backbone ViT-B/16 pré-entraîné avec DINO extrait 4 feature maps aux couches 3, 6, 9 et 12 via des forward hooks. L'image est découpée en 196 patches de 16$\times$16, projetés en embeddings de dimension 768. Le CLS token est retiré et les tokens reshapés en feature maps 2D : $(B, N, 768) \rightarrow (B, 768, 14, 14)$.

Tous les paramètres sont gelés, le mode \texttt{eval()} est forcé, et \texttt{@torch.no\_grad()} est appliqué au forward.

\subsection{Feature Pyramid Network (FPN)}

Deux FPN indépendants (RGB et Depth) transforment les features ViT via :
\begin{enumerate}
    \item \textbf{Connexions latérales} : Conv 1$\times$1 (768 $\to$ 256)
    \item \textbf{Pathway top-down} : Addition après interpolation bilinéaire
    \item \textbf{Lissage} : Conv 3$\times$3 + BatchNorm + ReLU
\end{enumerate}

\subsection{Cross-Modal Cross-Attention (CMCA)}

Le CMCA applique une attention croisée symétrique à chaque niveau FPN :
\begin{enumerate}
    \item Reshape spatial $\to$ séquence : $(B, C, H, W) \to (B, HW, C)$
    \item \textbf{2D$\to$3D} : $Q=\text{RGB}, K=V=\text{Depth}$ $\to$ RGB enrichi par Depth
    \item \textbf{3D$\to$2D} : $Q=\text{Depth}, K=V=\text{RGB}$ $\to$ Depth enrichi par RGB
    \item Résiduel + LayerNorm + FFN (256$\to$1024$\to$256, GELU) + Résiduel + LayerNorm
    \item Reshape séquence $\to$ spatial
\end{enumerate}

La formulation mathématique :
\begin{equation}
\text{RGB}' = \text{LN}\big(\text{RGB} + \text{MHA}(Q_{\text{RGB}}, K_{\text{Depth}}, V_{\text{Depth}})\big)
\end{equation}
\begin{equation}
\text{RGB}'' = \text{LN}\big(\text{RGB}' + \text{FFN}(\text{RGB}')\big)
\end{equation}

\subsection{Student Decoder (Reverse Distillation)}

Le Student Decoder reconstruit les features multi-échelles à partir de P4 (la plus profonde) :
\begin{itemize}
    \item \textbf{Bottleneck} : Conv 1$\times$1 + BN + ReLU
    \item \textbf{4 DecoderBlocks} : 2$\times$(Conv 3$\times$3 + BN + ReLU) + connexion résiduelle
    \item \textbf{Projections} : Conv 1$\times$1 + BN par niveau
    \item \textbf{Inversion} : [P4,P3,P2,P1] $\to$ [P1,P2,P3,P4]
\end{itemize}

\subsection{Adaptive Score Fusion}

Combinaison avec des poids appris par catégorie :
\begin{equation}
\mathbf{w} = \text{softmax}(\mathbf{W}[c]) \quad\Rightarrow\quad \text{fused} = w_{\text{rgb}} \cdot \text{map}_{\text{rgb}} + w_{\text{depth}} \cdot \text{map}_{\text{depth}}
\end{equation}

\subsection{Fonctions de loss}

La loss de Reverse Distillation combine deux composantes par modalité :
\begin{equation}
\mathcal{L}_{\text{modal}} = \underbrace{\frac{1}{L}\sum_{l=1}^{L}\left(1 - \cos(T_l, S_l)\right)}_{\text{Cosine Distance}} + \lambda_{L2} \cdot \underbrace{\frac{1}{L}\sum_{l=1}^{L} \|T_l - S_l\|^2}_{\text{MSE}}
\end{equation}
\begin{equation}
\mathcal{L}_{\text{total}} = w_{\text{rgb}} \cdot \mathcal{L}_{\text{rgb}} + w_{\text{depth}} \cdot \mathcal{L}_{\text{depth}}
\end{equation}
avec $\lambda_{L2} = 0.1$ et $w_{\text{rgb}} = w_{\text{depth}} = 1.0$.

\section{Architecture technique}

Le Tableau~\ref{tab:dimensions} trace les dimensions tensorielles à travers le réseau.

\begin{table}[H]
\centering
\caption{Dimensions tensorielles par module}
\label{tab:dimensions}
\begin{tabularx}{\textwidth}{X c}
\toprule
\textbf{Étape} & \textbf{Dimension} \\
\midrule
\rowcolor{tableheader} Entrée RGB / Depth & (B, 3, 224, 224) \\
Après ViT Hooks ($\times$4) & (B, 768, 14, 14) \\
\rowcolor{tableheader} Après FPN & (B, 256, 14, 14) \\
Après CMCA (séquence) & (B, 196, 256) \\
\rowcolor{tableheader} Après CMCA (spatial) & (B, 256, 14, 14) \\
Student Bottleneck & (B, 256, 14, 14) \\
\rowcolor{tableheader} Cosine Distance (par niveau) & (B, 1, 14, 14) \\
Upsample + Moyenne & (B, 1, 224, 224) \\
\rowcolor{tableheader} Anomaly Map (fusionnée) & (B, 1, 224, 224) \\
\bottomrule
\end{tabularx}
\end{table}

L'architecture logicielle est organisée en 4 couches : \textbf{Données} (dataset, géométrie), \textbf{Modèle} (7 modules assemblés par \texttt{cmad\_net.py}), \textbf{Entraînement} (pipeline 2 phases, loss), et \textbf{Évaluation} (10 métriques).

\section*{Conclusion du chapitre}
\addcontentsline{toc}{section}{Conclusion du chapitre}

Ce chapitre a présenté l'analyse et la conception complète de CMAD-Net. Les besoins fonctionnels et non-fonctionnels ont été formalisés. La conception détaille les 6 modules avec leurs spécifications, les diagrammes UML et les flux tensoriels. Le chapitre suivant présente la réalisation et les résultats.


% ============================================================
% CHAPITRE 4 : RÉALISATION
% ============================================================
\chapter{Réalisation}

\section*{Introduction du chapitre}
\addcontentsline{toc}{section}{Introduction du chapitre}

Ce dernier chapitre présente la réalisation concrète du projet. Nous détaillons les technologies, les interfaces du système, et les résultats expérimentaux sur MVTec 3D-AD.

\section{Technologies utilisées}

Le Tableau~\ref{tab:technologies} liste les technologies et bibliothèques utilisées.

\begin{table}[H]
\centering
\caption{Technologies et bibliothèques utilisées}
\label{tab:technologies}
\small
\begin{tabularx}{\textwidth}{l c X}
\toprule
\textbf{Technologie} & \textbf{Version} & \textbf{Rôle} \\
\midrule
\rowcolor{tableheader} Python & 3.10+ & Langage principal \\
PyTorch & $\geq$ 2.0 & Framework deep learning \\
\rowcolor{tableheader} torchvision & $\geq$ 0.15 & Transforms d'images \\
timm & $\geq$ 0.9.0 & Modèles ViT pré-entraînés (backbone) \\
\rowcolor{tableheader} NumPy & $\geq$ 1.24 & Calcul numérique \\
OpenCV & $\geq$ 4.7 & Traitement d'images (Sobel, morphologie) \\
\rowcolor{tableheader} scikit-learn & $\geq$ 1.3 & Métriques (AUROC, AP) \\
scikit-image & $\geq$ 0.19 & Region props pour AUPRO \\
\rowcolor{tableheader} SciPy & $\geq$ 1.10 & Gaussian filter \\
tifffile & $\geq$ 2023.2 & Lecture fichiers TIFF \\
\rowcolor{tableheader} Pandas & $\geq$ 2.0 & Données tabulaires \\
tqdm & $\geq$ 4.65 & Barres de progression \\
\rowcolor{tableheader} Matplotlib & $\geq$ 3.7 & Visualisation \\
\bottomrule
\end{tabularx}
\end{table}

Le Tableau~\ref{tab:technologies} montre un écosystème cohérent centré autour de PyTorch pour le deep learning.

\section{Réalisation}

\subsection{Configuration des hyperparamètres}

Le Tableau~\ref{tab:hyperparams} détaille les hyperparamètres.

\begin{table}[H]
\centering
\caption{Configuration des hyperparamètres}
\label{tab:hyperparams}
\begin{tabularx}{\textwidth}{X c c}
\toprule
\textbf{Hyperparamètre} & \textbf{Local} & \textbf{Kaggle} \\
\midrule
\rowcolor{tableheader} Image size & 224$\times$224 & 224$\times$224 \\
Batch size & 16 & 16 \\
\rowcolor{tableheader} Learning rate & 1e-4 & 1e-4 \\
Optimizer & AdamW & AdamW \\
\rowcolor{tableheader} Scheduler & CosineAnnealing & CosineAnnealing \\
Epochs (Phase 1) & 100 & 50 \\
\rowcolor{tableheader} Gradient clipping & max\_norm = 1.0 & max\_norm = 1.0 \\
FPN dimension & 256 & 256 \\
\rowcolor{tableheader} CMCA heads / dropout & 8 / 0.1 & 8 / 0.1 \\
$\lambda_{L2}$ (loss) & 0.1 & 0.1 \\
\rowcolor{tableheader} Gaussian $\sigma$ & 4.0 & 4.0 \\
Seed & 42 & 42 \\
\bottomrule
\end{tabularx}
\end{table}

\subsection{Sortie d'entraînement --- Phase 1 (Reverse Distillation)}

L'écran ci-dessous montre un extrait de la sortie de la Phase 1 :

\begin{lstlisting}[language={}, caption={Sortie d'entraînement Phase 1 (bagel)}, label={lst:phase1}]
CMAD-Net -- Training [1/10] : bagel
PHASE 1 : Entrainement Reverse Distillation
[Phase 1] Epoch  1/100, Loss: 0.552144
[Phase 1] Epoch  2/100, Loss: 0.126469
[Phase 1] Epoch  5/100, Loss: 0.020209
[Phase 1] Epoch 10/100, Loss: 0.004893
[Phase 1] Epoch 50/100, Loss: 0.000627
[Phase 1] Epoch100/100, Loss: 0.000439
\end{lstlisting}

Le Listing~\ref{lst:phase1} montre la convergence : réduction de 96\% en 5 epochs grâce au pré-entraînement ViT-DINO, puis plateau autour de 0.0004.

\subsection{Sortie de calibration --- Phase 2}

\begin{lstlisting}[language={}, caption={Sortie de calibration Phase 2}, label={lst:phase2}]
PHASE 2 : Calibration de la normalisation
Normalisation params: mean=0.1218, std=0.0520
\end{lstlisting}

Le Listing~\ref{lst:phase2} montre le calcul de $\mu$ et $\sigma$ sur le validation set pour la normalisation des cartes d'anomalies.

\subsection{Sortie d'évaluation}

\begin{lstlisting}[language={}, caption={Résultats pour bagel}, label={lst:eval}]
--- Results for bagel ---
 I-AUROC | P-AUROC | AUPRO@30% | I-AP  | I-F1
  0.936  |  0.969  |   0.792   | 0.984 | 0.935
\end{lstlisting}

Le Listing~\ref{lst:eval} montre les métriques pour la catégorie bagel : I-AUROC de 0.936 et P-AUROC de 0.969.

\subsection{Structure du code source}

Le projet totalise environ \textbf{2413 lignes} réparties en 15 fichiers Python :
\begin{itemize}
    \item \texttt{models/} : 7 fichiers, $\sim$943 lignes (architecture)
    \item \texttt{dataset/} : 3 fichiers, $\sim$392 lignes (données)
    \item \texttt{utils/} : 1 fichier, 129 lignes (loss)
    \item \texttt{evaluation/} : 3 fichiers, $\sim$575 lignes (métriques)
    \item \texttt{train\_cmad.py} : 374 lignes (script principal)
\end{itemize}

\section{Résultats expérimentaux}

\subsection{Résultats par catégorie}

Le Tableau~\ref{tab:resultats} présente les résultats sur les 10 catégories de MVTec 3D-AD.

\begin{table}[H]
\centering
\caption{Résultats expérimentaux par catégorie}
\label{tab:resultats}
\small
\begin{tabularx}{\textwidth}{X c c c c c c}
\toprule
\textbf{Catégorie} & \textbf{I-AUROC} & \textbf{P-AUROC} & \textbf{AUPRO@30\%} & \textbf{I-AP} & \textbf{I-F1} & \textbf{P-F1} \\
\midrule
\rowcolor{tableheader} bagel & 0.936 & 0.969 & 0.792 & 0.984 & 0.935 & 0.307 \\
cable\_gland & 0.612 & 0.939 & 0.807 & 0.864 & 0.895 & 0.085 \\
\rowcolor{tableheader} carrot & 0.750 & 0.977 & 0.922 & 0.919 & 0.917 & 0.152 \\
cookie & 0.885 & 0.965 & 0.872 & 0.961 & 0.922 & 0.331 \\
\rowcolor{tableheader} dowel & 0.757 & 0.962 & 0.847 & 0.921 & 0.900 & 0.100 \\
foam & 0.677 & 0.888 & 0.632 & 0.900 & 0.904 & 0.111 \\
\rowcolor{tableheader} peach & 0.735 & 0.959 & 0.842 & 0.895 & 0.897 & 0.107 \\
potato & 0.540 & 0.964 & 0.876 & 0.804 & 0.897 & 0.062 \\
\rowcolor{tableheader} rope & \textbf{0.953} & \textbf{0.982} & 0.866 & \textbf{0.983} & \textbf{0.936} & 0.278 \\
tire & 0.556 & 0.916 & 0.700 & 0.790 & 0.874 & 0.030 \\
\midrule
\textbf{Moyenne} & \textbf{0.740} & \textbf{0.952} & \textbf{0.816} & \textbf{0.902} & \textbf{0.908} & \textbf{0.156} \\
\bottomrule
\end{tabularx}
\end{table}

Le Tableau~\ref{tab:resultats} montre que CMAD-Net excelle en segmentation pixel (\textbf{P-AUROC $\geq$ 0.888} pour toutes les catégories). Les meilleures performances sont sur \textbf{rope} (I-AUROC=0.953) et \textbf{bagel} (0.936). Les catégories les plus difficiles sont \textbf{potato} (0.540) et \textbf{tire} (0.556).

\subsection{Moyennes globales}

Le Tableau~\ref{tab:moyennes} synthétise les performances moyennes.

\begin{table}[H]
\centering
\caption{Moyennes globales des métriques}
\label{tab:moyennes}
\begin{tabularx}{\textwidth}{X c c c}
\toprule
\textbf{Métrique} & \textbf{Moyenne} & \textbf{Meilleur} & \textbf{Pire} \\
\midrule
\rowcolor{tableheader} I-AUROC & 0.740 & rope (0.953) & potato (0.540) \\
\textbf{P-AUROC} & \textbf{0.952} & rope (0.982) & foam (0.888) \\
\rowcolor{tableheader} \textbf{AUPRO@30\%} & \textbf{0.816} & carrot (0.922) & foam (0.632) \\
AUPRO@10\% & 0.564 & carrot (0.765) & tire (0.379) \\
\rowcolor{tableheader} AUPRO@5\% & 0.382 & carrot (0.583) & tire (0.200) \\
AUPRO@1\% & 0.104 & carrot (0.181) & tire (0.035) \\
\rowcolor{tableheader} \textbf{I-AP} & \textbf{0.902} & bagel (0.984) & tire (0.790) \\
\textbf{I-F1} & \textbf{0.908} & rope (0.936) & tire (0.874) \\
\bottomrule
\end{tabularx}
\end{table}

Le Tableau~\ref{tab:moyennes} confirme les forces du modèle : \textbf{P-AUROC = 0.952} (excellente localisation pixel) et \textbf{I-AP = 0.902} (haute précision en classification).

\subsection{Paramètres de calibration}

Le Tableau~\ref{tab:calibration} présente les paramètres de normalisation calibrés.

\begin{table}[H]
\centering
\caption{Paramètres de calibration par catégorie}
\label{tab:calibration}
\begin{tabularx}{\textwidth}{X c c}
\toprule
\textbf{Catégorie} & $\boldsymbol{\mu}$ \textbf{(mean)} & $\boldsymbol{\sigma}$ \textbf{(std)} \\
\midrule
\rowcolor{tableheader} bagel & 0.1218 & 0.0520 \\
cable\_gland & 0.1440 & 0.0592 \\
\rowcolor{tableheader} carrot & 0.1224 & 0.0437 \\
cookie & 0.1326 & 0.0437 \\
\rowcolor{tableheader} dowel & 0.1227 & 0.0397 \\
foam & 0.1518 & 0.0397 \\
\rowcolor{tableheader} peach & 0.1153 & 0.0426 \\
potato & 0.1123 & 0.0399 \\
\rowcolor{tableheader} rope & 0.0988 & 0.0353 \\
tire & 0.1167 & 0.0380 \\
\bottomrule
\end{tabularx}
\end{table}

Le Tableau~\ref{tab:calibration} montre des valeurs moyennes homogènes (0.099--0.152), indiquant la stabilité du modèle.

\subsection{Analyse de la convergence}

La convergence se décompose en quatre phases :
\begin{enumerate}
    \item \textbf{Descente rapide} (epochs 1--5) : Loss de 0.55 à 0.02 (réduction de 96\%)
    \item \textbf{Convergence} (epochs 5--20) : Loss de 0.02 à 0.002
    \item \textbf{Ralentissement} (epochs 20--50) : Loss de 0.002 à 0.0006
    \item \textbf{Plateau} (epochs 50--100) : Loss stabilisée $\approx$ 0.0004
\end{enumerate}

La convergence rapide est caractéristique du paradigme Reverse Distillation avec backbone ViT pré-entraîné.

\subsection{Analyse par type de catégorie}

\begin{itemize}
    \item \textbf{Performants} (bagel, rope, cookie --- I-AUROC moyen 0.925) : anomalies bien contrastées
    \item \textbf{Intermédiaires} (carrot, dowel, peach --- I-AUROC moyen 0.747) : anomalies subtiles
    \item \textbf{Difficiles} (cable\_gland, foam, potato, tire --- I-AUROC moyen 0.596) : variabilité élevée
\end{itemize}

\section*{Conclusion du chapitre}
\addcontentsline{toc}{section}{Conclusion du chapitre}

Ce chapitre a présenté la réalisation de CMAD-Net. L'implémentation modulaire en $\sim$2400 lignes facilite la maintenance. Les résultats sur MVTec 3D-AD sont encourageants : \textbf{P-AUROC = 0.952}, \textbf{AUPRO@30\% = 0.816}, \textbf{I-AP = 0.902}. Les axes d'amélioration incluent l'entraînement GPU, l'augmentation de données et des backbones plus grands.


% ============================================================
% CONCLUSION GÉNÉRALE
% ============================================================
\chapter*{Conclusion Générale}
\addcontentsline{toc}{chapter}{Conclusion Générale}
\markboth{Conclusion Générale}{}

Ce projet de fin d'études a porté sur la conception, le développement et l'évaluation de \textbf{CMAD-Net} (Cross-Modal Attention Detection Network), une architecture pour la détection d'anomalies 3D multi-modale non-supervisée.

\subsection*{Bilan des réalisations}

Les principaux apports sont :
\begin{enumerate}
    \item \textbf{Architecture innovante} combinant ViT-B/16 DINOv2, FPN, CMCA, Student Decoder et fusion adaptative.
    \item \textbf{Fusion cross-modale profonde} via le CMCA (attention croisée symétrique).
    \item \textbf{Paradigme Reverse Distillation} pour une inférence rapide et une mémoire constante.
    \item \textbf{Résultats compétitifs} : P-AUROC = \textbf{0.952}, AUPRO@30\% = \textbf{0.816}.
    \item \textbf{Code modulaire} : $\sim$2400 lignes en modules indépendants.
\end{enumerate}

\subsection*{Objectifs atteints}

\begin{table}[H]
\centering
\begin{tabularx}{\textwidth}{c X c c}
\toprule
\textbf{N\textsuperscript{o}} & \textbf{Objectif} & \textbf{Statut} & \textbf{Résultat} \\
\midrule
\rowcolor{tableheader} O1 & Architecture multi-modale avec CMCA & \cmark & 4 niveaux $\times$ 8 têtes \\
O2 & Reverse Distillation & \cmark & Loss Cosine + MSE \\
\rowcolor{tableheader} O3 & Backbone ViT-B/16 DINOv2 & \cmark & Intégré via timm \\
O4 & FPN multi-échelles & \cmark & 4 niveaux top-down \\
\rowcolor{tableheader} O5 & Fusion adaptative & \cmark & Poids Softmax par catégorie \\
O6 & Évaluation 10 catégories & \cmark & Résultats complets \\
\rowcolor{tableheader} O7 & P-AUROC $>$ 0.95 & \cmark & \textbf{P-AUROC = 0.952} \\
\bottomrule
\end{tabularx}
\end{table}

\subsection*{Perspectives}

Pistes d'amélioration :
\begin{itemize}
    \item Backbone plus puissant (DINOv2-Large)
    \item Augmentation de données
    \item Cross-level attention dans le CMCA
    \item Test-Time Augmentation
    \item Déploiement temps réel (ONNX, TensorRT)
    \item Validation sur d'autres benchmarks
\end{itemize}

Ce projet démontre que la combinaison d'un Vision Transformer, d'un mécanisme d'attention croisée et du paradigme Reverse Distillation constitue une approche prometteuse pour la détection d'anomalies 3D multi-modale.


% ============================================================
% RÉFÉRENCES
% ============================================================
\begin{thebibliography}{99}
\addcontentsline{toc}{chapter}{Références}

\bibitem{dosovitskiy2021}
A. Dosovitskiy et al., ``An Image is Worth 16x16 Words: Transformers for Image Recognition at Scale'', \textit{ICLR}, 2021.

\bibitem{caron2021}
M. Caron et al., ``Emerging Properties in Self-Supervised Vision Transformers'', \textit{ICCV}, 2021.

\bibitem{oquab2023}
M. Oquab et al., ``DINOv2: Learning Robust Visual Features without Supervision'', \textit{arXiv:2304.07193}, 2023.

\bibitem{deng2022}
H. Deng et Y. Li, ``Anomaly Detection via Reverse Distillation from One-Class Embedding'', \textit{CVPR}, 2022.

\bibitem{lin2017}
T.-Y. Lin et al., ``Feature Pyramid Networks for Object Detection'', \textit{CVPR}, 2017.

\bibitem{bergmann2022}
P. Bergmann et al., ``The MVTec 3D-AD Dataset for Unsupervised 3D Anomaly Detection and Localization'', \textit{VISAPP}, 2022.

\bibitem{roth2022}
K. Roth et al., ``Towards Total Recall in Industrial Anomaly Detection'', \textit{CVPR}, 2022.

\bibitem{vaswani2017}
A. Vaswani et al., ``Attention Is All You Need'', \textit{NeurIPS}, 2017.

\bibitem{wang2021}
Y. Wang et al., ``Student-Teacher Feature Pyramid Matching for Anomaly Detection'', \textit{BMVC}, 2021.

\bibitem{bergmann2019}
P. Bergmann et al., ``MVTec AD -- A Comprehensive Real-World Dataset for Unsupervised Anomaly Detection'', \textit{CVPR}, 2019.

\bibitem{timm}
R. Wightman, ``PyTorch Image Models (timm)'', GitHub, 2019.

\bibitem{pytorch}
A. Paszke et al., ``PyTorch: An Imperative Style, High-Performance Deep Learning Library'', \textit{NeurIPS}, 2019.

\bibitem{adam}
D. P. Kingma et J. Ba, ``Adam: A Method for Stochastic Optimization'', \textit{ICLR}, 2015.

\bibitem{resnet}
K. He et al., ``Deep Residual Learning for Image Recognition'', \textit{CVPR}, 2016.

\bibitem{imagenet}
J. Deng et al., ``ImageNet: A Large-Scale Hierarchical Image Database'', \textit{CVPR}, 2009.

\end{thebibliography}


% ============================================================
% ANNEXES
% ============================================================
\appendix
\chapter{Hyperparamètres détaillés}

\begin{table}[H]
\centering
\caption{Configuration complète des hyperparamètres}
\begin{tabularx}{\textwidth}{X c X}
\toprule
\textbf{Hyperparamètre} & \textbf{Valeur} & \textbf{Justification} \\
\midrule
\rowcolor{tableheader} img\_size & 224 & Résolution native ViT-B/16 \\
fpn\_dim & 256 & Standard FPN \\
\rowcolor{tableheader} num\_heads & 8 & 8 têtes $\times$ 32 dim = 256 \\
dropout & 0.1 & Régularisation modérée \\
\rowcolor{tableheader} $\lambda_{L2}$ & 0.1 & MSE secondaire au cosine \\
$w_{\text{rgb}}, w_{\text{depth}}$ & 1.0, 1.0 & Poids égaux initiaux \\
\rowcolor{tableheader} lr & 1e-4 & Standard fine-tuning \\
weight\_decay & 1e-4 & Régularisation L2 \\
\rowcolor{tableheader} batch\_size & 16 & Limité par mémoire \\
gaussian\_sigma & 4.0 & Lissage pour 224$\times$224 \\
\rowcolor{tableheader} gradient\_clip & 1.0 & Stabilité \\
seed & 42 & Reproductibilité \\
\bottomrule
\end{tabularx}
\end{table}

\chapter{Glossaire}

\begin{longtable}{@{}p{4cm}p{10cm}@{}}
\toprule
\textbf{Terme} & \textbf{Définition} \\
\midrule
\endhead
AUROC & Aire sous la courbe ROC \\
AUPRO & Aire sous la courbe Per-Region Overlap \\
Backbone & Réseau pré-entraîné extracteur de features \\
Batch Normalization & Normalisation des activations par mini-batch \\
Cross-Attention & Attention où Q vient d'une source et K,V d'une autre \\
DINO & Auto-supervision par distillation sans labels \\
FPN & Feature Pyramid Network \\
Forward Hook & Callback PyTorch pour sorties intermédiaires \\
Ground Truth & Vérité terrain (annotations de référence) \\
Reverse Distillation & Student décodeur reconstruisant les features du teacher \\
Softmax & Normalisation convertissant scores en probabilités \\
Teacher-Student & Paradigme d'imitation de réseau \\
Vision Transformer & Transformer appliqué aux images en patches \\
\bottomrule
\end{longtable}


\end{document}
